\section{Running the flow}

\begin{frame}[fragile]{The golden path -- copy, paste, done}
  \begin{lstlisting}[style=shellstyle]
docker exec gf180 bash -lc '
  cd /foss/designs/template

  # One-time: clone the wafer-space PDK fork
  make clone-pdk

  # Every shell: activate GF180MCU
  source sak-pdk-script.sh gf180mcuD gf180mcu_fd_sc_mcu7t5v0

  # Run the full RTL-to-GDS flow on slot 1x1
  make librelane \
       SLOT=1x1 \
       PDK=gf180mcuD \
       PDK_ROOT=/foss/designs/template/gf180mcu
'
  \end{lstlisting}
  \textbf{Expected runtime:} $\sim$35--45\,min on a modern laptop.
  \note{%
    This is the entire flow. Three conceptual actions:
    (1) fetch the PDK fork, (2) activate the env, (3) invoke LibreLane
    through the template's Makefile. The \cmd{PDK\_ROOT} on the command
    line overrides whatever the Makefile default is -- deliberately
    redundant with \cmd{sak-pdk-script} to catch the trap we discuss in
    Section 6.\par
    Wall time breakdown on the audit run: Yosys 6\,s, floorplan 15\,s,
    PDN + placement + CTS + routing $\sim$3\,min, Magic DRC 32\,min,
    KLayout antenna 7\,min, LVS + STA + IR 1\,min. If anyone asks
    ``why so slow'': Magic DRC on 84k standard cells with padring. The
    \cmd{make librelane-nodrc} target drops DRC for iteration.}
\end{frame}

\begin{frame}{What scrolls by while it runs}
  \begin{itemize}\small
    \item LibreLane prints a banner and numbered steps
          (\cmd{Step 1: Yosys.Synthesis}, \ldots).
    \item Yosys: inferred cell counts + area.
    \item OpenROAD: floorplan, PDN, placement, CTS, routing, STA.
    \item Magic + KLayout: DRC (the slow parts).
    \item Netgen: LVS.
    \item Final banner: run directory + pointer to \cmd{metrics.csv}.
  \end{itemize}
  \vspace{0.4em}
  \screenshotplaceholder{slides/figures/librelane_run_output.png}{%
    Tail of a successful \cmd{make librelane} run.}
  \note{%
    The log is dense but readable. Every step prints a header like
    \cmd{== Step 42: OpenROAD.Floorplan} so you can grep. Warnings during
    routing repair iterations are expected; failures manifest as the flow
    stopping abruptly on a step boundary.\par
    The final banner tells you two things: the timestamped run directory
    (where the full logs and intermediate DEFs live) and the path to
    \cmd{metrics.csv}. Bookmark both when presenting live.}
\end{frame}

\begin{frame}[fragile]{Where artifacts land}
  \begin{lstlisting}[basicstyle=\ttfamily\footnotesize,frame=single,
                     backgroundcolor=\color{softbg}]
template/
  librelane/runs/RUN_<timestamp>/   # per-run: logs, DEFs, reports
  final/                            # merged: GDS, LEF, SDC, SDF, SPEF,
                                    # metrics.csv, render/chip_top.png
  \end{lstlisting}
  Useful after a successful run:
  \begin{itemize}\small
    \item \cmd{final/gds/chip\_top.gds} -- the deliverable.
    \item \cmd{final/metrics.csv} -- one row per metric.
    \item \cmd{final/render/chip\_top.png} -- auto-render for posters.
  \end{itemize}
  \note{%
    The split matters: \cmd{runs/} is every flow artifact organised per
    run (useful for bisecting), \cmd{final/} is the curated deliverable
    set that LibreLane copies via its \cmd{--save-views-to} flag. You almost
    never dig into \cmd{runs/} unless something failed.\par
    Practical tip: \cmd{runs/} balloons fast. Add it to \cmd{.gitignore}.
    Commit only a trimmed \cmd{final/metrics.csv} as a regression baseline.
    Add a CI job that greps the next run's \cmd{metrics.csv} for the key
    fields and compares -- that is how you catch silent regressions.}
\end{frame}

\begin{frame}{79 steps grouped into six phases}
  \begin{center}
  \begin{tikzpicture}[node distance=0.35cm and 0.3cm]
    \node[pipelinebox,  minimum width=2.3cm] (s1) {1. Synthesis\\\scriptsize Yosys};
    \node[pipelinebox,  right=of s1, minimum width=2.3cm] (s2) {2. Floorplan\\\scriptsize OpenROAD};
    \node[pipelinebox,  right=of s2, minimum width=2.3cm] (s3) {3. PDN\\\scriptsize OpenROAD};
    \node[pipelinebox,  below=of s3, minimum width=2.3cm] (s4) {4. Place \& CTS\\\scriptsize OpenROAD};
    \node[pipelinebox,  left=of s4, minimum width=2.3cm] (s5) {5. Routing\\\scriptsize OpenROAD};
    \node[pipelinebox3, left=of s5, minimum width=2.3cm] (s6) {6. Signoff\\\scriptsize Magic / KLayout / Netgen};
    \draw[arrow] (s1) -- (s2);
    \draw[arrow] (s2) -- (s3);
    \draw[arrow] (s3) -- (s4);
    \draw[arrow] (s4) -- (s5);
    \draw[arrow] (s5) -- (s6);
  \end{tikzpicture}
  \end{center}
  \vspace{0.3em}
  Each phase writes a milestone file. Any failure halts the flow.
  \note{%
    This is the recap diagram before the deep-dive section. Audience
    already saw it on slide 2 -- the difference now is that they know
    LibreLane is the driver and they have seen the command that starts it.
    Point at each box and promise: ``the next five slides tell you what
    actually happens inside these.''\par
    If pressed: why 79 steps and not 6? LibreLane breaks the phases into
    many auxiliary passes (netlist cleanup, antenna repair, resizing,
    multiple DRC runs). The 6-phase model is the mental pipeline; the 79
    steps are the implementation.}
\end{frame}
